```latex
\documentclass{article}

% Language setting
% Replace `english' with e.g. `spanish' to change the document language
\usepackage[portuguese]{babel}

% Set page size and margins
% Replace `letterpaper' with `a4paper' for UK/EU standard size
\usepackage[letterpaper,top=2cm,bottom=2cm,left=3cm,right=3cm,marginparwidth=1.75cm]{geometry}

% Useful packages
\usepackage{amsmath}
\usepackage{graphicx}
\usepackage[colorlinks=true, allcolors=blue]{hyperref}
\usepackage{float}

\title{Métodos Numéricos\\Trabalho II}
\author{Pedro Fam Haggstram\thanks{\texttt{pedro.haggstram@edu.pucrs.br}}, Antônio Pasquali Coelho\thanks{\texttt{antonio.pasquali@edu.pucrs.br}}\\
Escola Politécnica — PUCRS}

\begin{document}
\maketitle

\begin{abstract}
Este relatório descreve a solução obtida exclusivamente por meio de computação simbólica em \textit{Common Lisp} para calcular a altura de um chapéu lançado ao ar no instante $t = 1{,}3$ segundos. A partir dos pontos experimentais fornecidos por Mariela, construímos o polinômio interpolador de Lagrange em linguagem Lisp, geramos seu formato compatível com \texttt{gnuplot} e, em seguida, plotamos o gráfico resultante. Por fim, discutimos a expressão simbólica, identificamos máximos locais e globais do polinômio e avaliamos $f(1{,}3)$. Também abordamos o conceito de \emph{overfitting}, que surge quando polinômios de alto grau oscilam demais e capturam ruído nos dados.
\end{abstract}

\section{Introdução}
Mariela anotou as estimativas de altura do seu chapéu em diferentes instantes após o lançamento. Os dados são:
\[
\begin{aligned}
f(0\,\text{s}) &= 1{,}6\,\text{m},\\
f(0{,}5\,\text{s}) &= 3\,\text{m},\\
f(0{,}8\,\text{s}) &= 3{,}7\,\text{m},\\
f(1\,\text{s})   &= 4{,}8\,\text{m},\\
f(1{,}4\,\text{s}) &= 4{,}7\,\text{m},\\
f(2\,\text{s})   &= 3\,\text{m},\\
f(2{,}7\,\text{s}) &= 2{,}6\,\text{m}.
\end{aligned}
\]
O objetivo é usar interpolação de Lagrange para estimar $f(1{,}3)$ e verificar se o máximo ocorre de fato em torno desse instante. Toda a implementação foi realizada em Lisp, gerando uma expressão simbólica de grau 6 que depois transformamos numa string própria para \texttt{gnuplot} e produzimos o gráfico final.

\section{Interação com ChatGPT}
Para comparar estratégias, obtivemos primeiro um valor de referência de $f(1{,}3)$ enviando este prompt ao ChatGPT:
\begin{quote}
\textit{“Mariela jogou seu chapéu para cima e anotou:
\[
\begin{aligned}
f(0)&=1{,}6,\quad f(0{,}5)=3,\quad f(0{,}8)=3{,}7,\quad f(1)=4{,}8,\\
f(1{,}4)&=4{,}7,\quad f(2)=3,\quad f(2{,}7)=2{,}6.
\end{aligned}
\]
Ela acredita que o chapéu está no ponto mais alto em $t=1{,}3$\,s. Use interpolação de Lagrange ou Newton para calcular $f(1{,}3)$. Também responda se esses pontos correspondem a uma parábola.”}
\end{quote}

O ChatGPT explicou a fórmula de Lagrange, alertou sobre o Fenômeno de Runge ao usar todos os sete pontos (polinômio de grau 6) e sugeriu, para maior estabilidade local, usar um subconjunto de quatro pontos em torno de $t=1{,}3$. Com esse método cúbico, obteve
\[
f_{\mathrm{ChatGPT}}(1{,}3) \approx 4{,}93\;\text{m}.
\]
Quanto à pergunta sobre parábola, o modelo apontou que, embora o movimento ideal de projétil seja quadrático, os pontos de Mariela não se ajustam perfeitamente a um único polinômio de grau 2 devido a pequenas imprecisões. Sugeriu que, em geral, os dados se aproximam de uma parábola cujo vértice ficaria entre $t \approx 1{,}1$ e $1{,}3$\,s.

\section{Implementação em Lisp (Computação Simbólica)}
Nesta seção, mostramos como geramos o polinômio de Lagrange simbolicamente em \textit{Common Lisp}, transformamos a expressão num formato legível pelo \texttt{gnuplot} e, então, avaliamos $f(1{,}3)$ e extraímos características do polinômio (máximos locais e globais).

\subsection{Geração simbólica do polinômio}
No Lisp, definimos a função a seguir para construir o interpolador de Lagrange a partir da lista de pares $(x_i, y_i)$:
\begin{verbatim}
(defun lagrange-interpolating-polynomial (points)
  "Retorna uma expressão simbólica do polinômio de Lagrange para os pontos fornecidos."
  (let ((terms
         (loop for i from 0 below (length points)
               for (xi yi) = (nth i points)
               collect
               `(* ,yi
                   ,(reduce (lambda (a b) `(* ,a ,b))
                            (loop for j from 0 below (length points)
                                  for (xj _) = (nth j points)
                                  when (/= i j)
                                  collect
                                  `(/ (- x ,xj) (- ,xi ,xj))))))))
    `(+ ,@terms)))
\end{verbatim}

Em seguida, definimos os pontos de Mariela e chamamos a função:
\begin{verbatim}
(defparameter pontos
  '((0   1.6)
    (0.5 3  )
    (0.8 3.7)
    (1   4.8)
    (1.4 4.7)
    (2   3  )
    (2.7 2.6)))

(defparameter poly
  (lagrange-interpolating-polynomial pontos))
\end{verbatim}

O resultado \texttt{poly} é uma lista de Lisp que representa, simbolicamente, o polinômio de grau 6:
\[
f_{\mathrm{Lisp}}(x) \;=\; \sum_{i=0}^{6} y_i \,\prod_{\substack{0 \leq j \leq 6 \\ j \neq i}} \frac{x - x_j}{x_i - x_j}.
\]

\subsection{Conversão para formato \texttt{gnuplot}}
Para plotagem, precisamos de uma string em que as operações estejam escritas em notação familiar ao \texttt{gnuplot}. Criamos a função \texttt{p-print-gnuplot} que percorre a expressão simbólica e monta uma string correspondente. O gráfico resultante está na Figura~\ref{fig:graph1}.

\subsection{Avaliação simbólica em $x = 1{,}3$}
Para extrair numericamente $f(1{,}3)$, definimos uma função que substitui $x$ por $1{,}3$ na expressão e avalia:
\begin{verbatim}
(defun eval-at-x (expr x-value)
  "Avalia a expressão simbólica EXPR após substituir x por x-value."
  (eval (subst x-value 'x expr)))

(format t "f(1.3) = ~a~%" (eval-at-x poly 1.3))
\end{verbatim}

O valor resultante foi:
\[
f_{\mathrm{Lisp}}(1{,}3) \approx 5{,}130037\;\text{m},
\]
confirmando que o polinômio de grau 6 (usando todos os pontos) gera um valor levemente diferente daquele de grau 3 extraído do ChatGPT.

\subsection{Propriedades do polinômio de grau 6}
Usando ferramentas gráficas (por exemplo, métodos de busca numérica ou visualização em \texttt{gnuplot}), encontramos que o polinômio gerado possui:
\begin{itemize}
    \item \textbf{Máximo local} em $x \approx 1{,}19388$ com $y \approx 5{,}29248$, valor este ainda maior que $f(1{,}3)$.
    \item \textbf{Máximo global} em $\bigl(x \approx 2{,}48358,\;y \approx 9{,}89273\bigr)$, muito acima dos dados iniciais. 
\end{itemize}
Como o polinômio é de grau 6, ele não segue a forma de uma parábola (que seria esperada para um projétil ideal). Polinômios de grau elevado podem oscilar entre os pontos de interpolação, fenômeno que se intensifica quando os dados possuem ruído ou imprecisões. Nesses casos, dizemos que a interpolação pode estar sofrendo de \emph{overfitting}.

\subsubsection*{O que é \emph{overfitting}?}
\emph{Overfitting} ocorre quando um modelo (neste caso, um polinômio de alto grau) se ajusta não apenas ao comportamento geral dos dados, mas também a ruídos ou variações locais que não representam a tendência verdadeira. Como resultado, o polinômio pode apresentar oscilações fora do intervalo dos pontos observados, gerando máximos ou mínimos espúrios. Em contextos de interpolação, o \emph{overfitting} torna o resultado menos confiável para predição em novos valores de $x$, pois o modelo ficou “sobreajustado” aos detalhes idiossincráticos dos pontos fornecidos.

\subsection{Plotagem no \texttt{gnuplot}}
Com a string gerada por \texttt{p-print-gnuplot}, criamos um arquivo texto contendo:
\begin{verbatim}
f(x) = <conteúdo de string-gnuplot>
set terminal jpeg size 800,600
set output "our_graph.jpg"
set xlabel "Tempo (s)"
set ylabel "Altura (m)"
set title "Interpolação de Lagrange (grau 6) em Lisp"
plot f(x) with lines lw 2
\end{verbatim}
Ao executar no \texttt{gnuplot}, produzimos a Figura~\ref{fig:graph1}, que mostra a curva interpoladora completa em função dos sete pontos originais.

\begin{figure}[htp]
    \centering
    \includegraphics[width=12cm]{our_graph.jpg}
    \caption{Gráfico da interpolação de Lagrange (grau 6) gerado a partir da expressão simbólica produzida em Lisp.}
    \label{fig:graph1}
\end{figure}

\begin{figure}[htp]
    \centering
    \includegraphics[width=12cm]{gpt_graph.jpg}
    \caption{Gráfico do ChatGPT (interpolação cúbica com subconjunto de 4 pontos).}
    \label{fig:graph2}
\end{figure}

\section{Comparação e conclusões}
Ao comparar os valores obtidos:
\begin{itemize}
    \item \textbf{ChatGPT (grau 3, 4 pontos)}: $f(1{,}3) \approx 4{,}93$\,m.
    \item \textbf{Lisp (grau 6, 7 pontos)}: $f(1{,}3) \approx 5{,}13$\,m.
\end{itemize}
Além disso, o polinômio de grau 6 exibe oscilações que produzem um máximo local em $x \approx 1{,}19388$ com $y \approx 5{,}29248$ e um máximo global em $\bigl(x \approx 2{,}48358,\;y \approx 9{,}89273\bigr)$. Tais comportamentos não são esperados num movimento parabólico real de um chapéu, mas surgem devido ao grau elevado do interpolador e às pequenas imprecisões nos pontos originais.

\emph{Overfitting} explica por que o polinômio de grau 6 gere picos que não correspondem fisicamente à trajetória do chapéu: o modelo polinomial ficou excessivamente ajustado aos detalhes dos dados, capturando ruído que introduz oscilações artefatuosas. Por isso, ao estimar $f(1{,}3)$, a interpolação cúbica local (grau 3) demonstra-se mais estável e confiável.

Concluímos que, embora a interpolação polinomial complete seja viável simbolicamente em Lisp, deve-se ter cautela com polinômios de grau alto em dados experimentais, pois o \emph{overfitting} pode gerar máximos espúrios e resultados pouco realistas. A escolha de um grau menor ou de métodos de ajuste por mínimos quadrados pode produzir estimativas mais consistentes com a física do movimento.

\end{document}
```
